\begin{register}{H}{LP\_CLKRST\_LP\_CLK\_CONF\_REG}{0x0000}\label{regdesc:LPCLKRSTLPCLKCONFREG}
\regfield{(reserved)}{20}{4}{0000000000000000000000000000}%
\regfield{LP\_CLKRST\_FAST\_CLK\_SEL}{2}{2}{{0x1}}%
\regfield{LP\_CLKRST\_SLOW\_CLK\_SEL}{2}{0}{{0x0}}%
\reglabel{Reset}\regnewline%
\begin{regdesc}\begin{reglist}
\label{fielddesc:LPCLKRSTSLOWCLKSEL}\item [LP\_CLKRST\_SLOW\_CLK\_SEL] 配置 LP\_SLOW\_CLK 的时钟源。\\
0：RC\_SLOW\_CLK \\
1：XTAL32K\_CLK \\
2：OSC\_SLOW\_CLK  \\
3：无效值。没有作用 \\
 (R/W)
\label{fielddesc:LPCLKRSTFASTCLKSEL}\item [LP\_CLKRST\_FAST\_CLK\_SEL] 配置 LP\_FAST\_CLK 的时钟源。\\
0：RC\_FAST\_CLK \\
1：XTAL\_D2\_CLK  \\ (R/W)
\end{reglist}\end{regdesc}
\end{register}


\begin{register}{H}{LP\_CLKRST\_LP\_CLK\_PO\_EN\_REG}{0x0004}\label{regdesc:LPCLKRSTLPCLKPOENREG}
\regfield{(reserved)}{21}{11}{000000000000000000000}%
\regfield{LP\_CLKRST\_LPBUS\_OEN}{1}{10}{{1}}%
\regfield{LP\_CLKRST\_RNG\_OEN}{1}{9}{{1}}%
\regfield{LP\_CLKRST\_FAST\_OEN}{1}{8}{{1}}%
\regfield{LP\_CLKRST\_SLOW\_OEN}{1}{7}{{1}}%
\regfield{LP\_CLKRST\_CORE\_EFUSE\_OEN}{1}{6}{{1}}%
\regfield{LP\_CLKRST\_XTAL32K\_OEN}{1}{5}{{1}}%
\regfield{(reserved))}{1}{4}{{1}}%
\regfield{LP\_CLKRST\_FOSC\_OEN}{1}{3}{{1}}%
\regfield{LP\_CLKRST\_SOSC\_OEN}{1}{2}{{1}}%
\regfield{LP\_CLKRST\_AON\_FAST\_OEN}{1}{1}{{1}}%
\regfield{LP\_CLKRST\_AON\_SLOW\_OEN}{1}{0}{{1}}%
\reglabel{Reset}\regnewline%
\begin{regdesc}\begin{reglist}
\label{fielddesc:LPCLKRSTAONSLOWOEN}\item [LP\_CLKRST\_AON\_SLOW\_OEN] 配置是否使能 LP\_DYN\_SLOW\_CLK 输出到 pad 的时钟门控。\\
0：禁止时钟信号通过\\
1：使能时钟信号通过 \\ (R/W)
\label{fielddesc:LPCLKRSTAONFASTOEN}\item [LP\_CLKRST\_AON\_FAST\_OEN] 配置是否使能 LP\_DYN\_FAST\_CLK 输出到 pad 的时钟门控。\\
0：禁止时钟信号通过\\
1：使能时钟信号通过 \\ (R/W)
\label{fielddesc:LPCLKRSTSOSCOEN}\item [LP\_CLKRST\_SOSC\_OEN] 配置是否使能 OSC\_SLOW\_CLK 输出到 pad 的时钟门控。\\
0：禁止时钟信号通过\\
1：使能时钟信号通过 \\ (R/W)
\label{fielddesc:LPCLKRSTFOSCOEN}\item [LP\_CLKRST\_FOSC\_OEN] 配置是否使能 RC\_FAST\_CLK 输出到 pad 的时钟门控。\\
0：禁止时钟信号通过\\
1：使能时钟信号通过 \\ (R/W)
\label{fielddesc:LPCLKRSTXTAL32KOEN}\item [LP\_CLKRST\_XTAL32K\_OEN] 配置是否使能 XTAL32K\_CLK 输出到 pad 的时钟门控。\\
0：禁止时钟信号通过\\
1：使能时钟信号通过 \\ (R/W)

\item [见下页...]
\end{reglist}\end{regdesc}
\end{register}

\addtocounter{Regfloat}{-1}
\begin{register}{H}{LP\_CLKRST\_LP\_CLK\_PO\_EN\_REG}{0x0004}
\begin{regdesc}\begin{reglist}
\item [接上页...]

\label{fielddesc:LPCLKRSTCOREEFUSEOEN}\item [LP\_CLKRST\_CORE\_EFUSE\_OEN] 配置是否使能 EFUSE\_CTRL 时钟输出到 pad 的时钟门控。\\
0：禁止时钟信号通过\\
1：使能时钟信号通过 \\ (R/W)
\label{fielddesc:LPCLKRSTSLOWOEN}\item [LP\_CLKRST\_SLOW\_OEN] 配置是否使能 LP\_SLOW\_CLK 输出到 pad 的时钟门控。\\
0：禁止时钟信号通过\\
1：使能时钟信号通过 \\ (R/W)
\label{fielddesc:LPCLKRSTFASTOEN}\item [LP\_CLKRST\_FAST\_OEN] 配置是否使能 LP\_FAST\_CLK 输出到 pad 的时钟门控。\\
0：禁止时钟信号通过\\
1：使能时钟信号通过 \\ (R/W)
\label{fielddesc:LPCLKRSTRNGOEN}\item [LP\_CLKRST\_RNG\_OEN] 配置是否使能 RNG 时钟输出到 pad 的时钟门控。\\
0：禁止时钟信号通过\\
1：使能时钟信号通过 \\ (R/W)
\label{fielddesc:LPCLKRSTLPBUSOEN}\item [LP\_CLKRST\_LPBUS\_OEN] 配置是否使能 LP 总线时钟输出到 pad 的时钟门控。\\
0：禁止时钟信号通过\\
1：使能时钟信号通过 \\ (R/W)
\end{reglist}\end{regdesc}
\end{register}


\begin{register}{H}{LP\_CLKRST\_LP\_CLK\_EN\_REG}{0x0008}\label{regdesc:LPCLKRSTLPCLKENREG}
\regfield{LP\_CLKRST\_FAST\_ORI\_GATE}{1}{31}{{0}}%
\regfield{(reserved)}{31}{0}{0000000000000000000000000000000}%
\reglabel{Reset}\regnewline%
\begin{regdesc}\begin{reglist}
\label{fielddesc:LPCLKRSTFASTORIGATE}\item [LP\_CLKRST\_FAST\_ORI\_GATE] 配置 LP\_FAST\_CLK 的时钟门控。\\
0：无效。时钟门控由硬件 FSM 控制\\
1：强制时钟通过门控 \\ (R/W)
\end{reglist}\end{regdesc}
\end{register}


\begin{register}{H}{LP\_CLKRST\_LP\_RST\_EN\_REG}{0x000C}\label{regdesc:LPCLKRSTLPRSTENREG}
\regfield{LP\_CLKRST\_ANA\_PERI\_RESET\_EN}{1}{31}{{0}}%
\regfield{LP\_CLKRST\_WDT\_RESET\_EN}{1}{30}{{0}}%
\regfield{LP\_CLKRST\_LP\_TIMER\_RESET\_EN}{1}{29}{{0}}%
\regfield{LP\_CLKRST\_AON\_EFUSE\_CORE\_RESET\_EN}{1}{28}{{0}}%
\regfield{(reserved)}{28}{0}{0000000000000000000000000000}%
\reglabel{Reset}\regnewline%
\begin{regdesc}\begin{reglist}
\label{fielddesc:LPCLKRSTAONEFUSECORERESETEN}\item [LP\_CLKRST\_AON\_EFUSE\_CORE\_RESET\_EN] 配置复位 EFUSE\_CTRL 常开部分。\\
0：无效 \\
1：复位  \\ (R/W)
\label{fielddesc:LPCLKRSTLPTIMERRESETEN}\item [LP\_CLKRST\_LP\_TIMER\_RESET\_EN] 配置复位 RTC 定时器。\\
0：无效 \\
1：复位  \\ (R/W)
\label{fielddesc:LPCLKRSTWDTRESETEN}\item [LP\_CLKRST\_WDT\_RESET\_EN] 配置复位 RWDT 和超级看门狗定时器。\\
0：无效 \\
1：复位  \\ (R/W)
\label{fielddesc:LPCLKRSTANAPERIRESETEN}\item [LP\_CLKRST\_ANA\_PERI\_RESET\_EN] 配置复位模拟外设，包括欠压检测器。\\
0：无效 \\
1：复位  \\ (R/W)
\end{reglist}\end{regdesc}
\end{register}


\begin{register}{H}{LP\_CLKRST\_RESET\_CAUSE\_REG}{0x0010}\label{regdesc:LPCLKRSTRESETCAUSEREG}
\regfield{LP\_CLKRST\_CORE0\_RESET\_FLAG\_CLR}{1}{31}{{0}}%
\regfield{LP\_CLKRST\_CORE0\_RESET\_FLAG\_SET}{1}{30}{{0}}%
\regfield{LP\_CLKRST\_CORE0\_RESET\_CAUSE\_CLR}{1}{29}{{0}}%
\regfield{(reserved)}{23}{6}{00000000000000000000000}%
\regfield{LP\_CLKRST\_CORE0\_RESET\_FLAG}{1}{5}{{1}}%
\regfield{LP\_CLKRST\_RESET\_CAUSE}{5}{0}{{0}}%
\reglabel{Reset}\regnewline%
\begin{regdesc}\begin{reglist}
\label{fielddesc:RTCCLKRSTRESETCAUSE}\item [RTC\_CLKRST\_RESET\_CAUSE] 表示复位原因。(RO)
\label{fielddesc:LPCLKRSTCORE0RESETFLAG}\item [LP\_CLKRST\_CORE0\_RESET\_FLAG] 表示发生的复位是否为 LP\_CLKRST\_RESET\_CAUSE 中记录的复位。\\
0：非法复位\\
1：已记录的复位\\ (RO)
\label{fielddesc:LPCLKRSTCORE0RESETCAUSECLR}\item [LP\_CLKRST\_CORE0\_RESET\_CAUSE\_CLR] 写 1 清除 LP\_CLKRST\_RESET\_CAUSE。 (WT)
\label{fielddesc:LPCLKRSTCORE0RESETFLAGSET}\item [LP\_CLKRST\_CORE0\_RESET\_FLAG\_SET] 写 1 置位
 LP\_CLKRST\_CORE0\_RESET\_FLAG。 (WT)
\label{fielddesc:LPCLKRSTCORE0RESETFLAGCLR}\item [LP\_CLKRST\_CORE0\_RESET\_FLAG\_CLR] 写 1 清除
 LP\_CLKRST\_CORE0\_RESET\_FLAG。 (WT)
\end{reglist}\end{regdesc}
\end{register}


\begin{register}{H}{LP\_CLKRST\_CPU\_RESET\_REG}{0x0014}\label{regdesc:LPCLKRSTCPURESETREG}
\regfield{LP\_CLKRST\_CPU\_STALL\_EN}{1}{31}{{0}}%
\regfield{LP\_CLKRST\_CPU\_STALL\_WAIT}{5}{26}{{1}}%
\regfield{LP\_CLKRST\_RTC\_WDT\_CPU\_RESET\_EN}{1}{25}{{0}}%
\regfield{LP\_CLKRST\_RTC\_WDT\_CPU\_RESET\_LENGTH}{3}{22}{{1}}%
\regfield{(reserved)}{22}{0}{0000000000000000000000}%
\reglabel{Reset}\regnewline%
\begin{regdesc}\begin{reglist}
\label{fielddesc:LPCLKRSTRTCWDTCPURESETLENGTH}\item [LP\_CLKRST\_RTC\_WDT\_CPU\_RESET\_LENGTH] 配置 RWDT 复位 CPU 的时间长度。\\
单位：LP\_DYN\_FAST\_CLK 周期。 \\ (R/W)
\label{fielddesc:LPCLKRSTRTCWDTCPURESETEN}\item [LP\_CLKRST\_RTC\_WDT\_CPU\_RESET\_EN] 配置 RWDT 是否能复位 CPU。\\
0：RWDT 超时时无法复位 CPU \\
1：RWDT 超时时复位 CPU\\ (R/W)
\label{fielddesc:LPCLKRSTCPUSTALLWAIT}\item [LP\_CLKRST\_CPU\_STALL\_WAIT] 配置 CPU 停顿状态 (CPU Stall) 到复位之间的时间间隔。\\
单位：LP\_DYN\_FAST\_CLK 周期。 \\ (R/W)
\label{fielddesc:LPCLKRSTCPUSTALLEN}\item [LP\_CLKRST\_CPU\_STALL\_EN] 配置在 RWDT 和软件复位 CPU 之前是否进入 CPU 停顿状态。\\
0：否 \\
1：是  \\ (R/W)
\end{reglist}\end{regdesc}
\end{register}


\begin{register}{H}{LP\_CLKRST\_FOSC\_CNTL\_REG}{0x0018}\label{regdesc:LPCLKRSTFOSCCNTLREG}
\regfield{LP\_CLKRST\_FOSC\_DFREQ}{10}{22}{{600}}%
\regfield{(reserved)}{22}{0}{0000000000000000000000}%
\reglabel{Reset}\regnewline%
\begin{regdesc}\begin{reglist}
\label{fielddesc:LPCLKRSTFOSCDFREQ}\item [LP\_CLKRST\_FOSC\_DFREQ] 配置 RC\_FAST\_CLK 频率，值越大，时钟周期越大。(R/W)
\end{reglist}\end{regdesc}
\end{register}




\begin{register}{H}{LP\_CLKRST\_CLK\_TO\_HP\_REG}{0x0020}\label{regdesc:LPCLKRSTCLKTOHPREG}
\regfield{LP\_CLKRST\_ICG\_HP\_FOSC}{1}{31}{{1}}%
\regfield{(reserved))}{1}{30}{{1}}%
\regfield{LP\_CLKRST\_ICG\_HP\_SOSC}{1}{29}{{1}}%
\regfield{LP\_CLKRST\_ICG\_HP\_XTAL32K}{1}{28}{{1}}%
\regfield{(reserved)}{28}{0}{0000000000000000000000000000}%
\reglabel{Reset}\regnewline%
\begin{regdesc}\begin{reglist}
\label{fielddesc:LPCLKRSTICGHPXTAL32K}\item [LP\_CLKRST\_ICG\_HP\_XTAL32K] 配置 XTAL32K\_CLK 到 HP System 的时钟门控。\\
0：时钟无法到达 HP System \\
1：时钟能够到达 HP System  \\ (R/W)
\label{fielddesc:LPCLKRSTICGHPSOSC}\item [LP\_CLKRST\_ICG\_HP\_SOSC] 配置 RC\_SLOW\_CLK 到 HP System 的时钟门控。\\
0：时钟无法到达 HP System \\
1：时钟能够到达 HP System  \\ (R/W)
\label{fielddesc:LPCLKRSTICGHPFOSC}\item [LP\_CLKRST\_ICG\_HP\_FOSC] 配置 RC\_FAST\_CLK 到 HP System 的时钟门控。\\
0：时钟无法到达 HP System \\
1：时钟能够到达 HP System \\ (R/W)
\end{reglist}\end{regdesc}
\end{register}


\begin{register}{H}{LP\_CLKRST\_LPMEM\_FORCE\_REG}{0x0024}\label{regdesc:LPCLKRSTLPMEMFORCEREG}
\regfield{LP\_CLKRST\_LPMEM\_CLK\_FORCE\_ON}{1}{31}{{0}}%
\regfield{(reserved)}{31}{0}{0000000000000000000000000000000}%
\reglabel{Reset}\regnewline%
\begin{regdesc}\begin{reglist}
\label{fielddesc:LPCLKRSTLPMEMCLKFORCEON}\item [LP\_CLKRST\_LPMEM\_CLK\_FORCE\_ON] 配置强制打开 LP 存储器的时钟门控。\\
0：无效。时钟门控由硬件 FSM 控制\\
1：强制打开时钟门控  \\ (R/W)
\end{reglist}\end{regdesc}
\end{register}

\begin{register}{H}{LP\_CLKRST\_LPPERI\_REG}{0x0028}\label{regdesc:LPCLKRSTLPPERIREG}
\regfield{(reserved)}{2}{30}{{0}}%
\regfield{LP\_CLKRST\_LP\_SEL\_XTAL32K}{1}{29}{{1}}%
\regfield{LP\_CLKRST\_LP\_SEL\_XTAL}{1}{28}{{0}}%
\regfield{LP\_CLKRST\_LP\_SEL\_OSC\_FAST}{1}{27}{{0}}%
\regfield{LP\_CLKRST\_LP\_SEL\_OSC\_SLOW}{1}{26}{{0}}%
\regfield{(reserved)}{2}{24}{{0}}%
\regfield{LP\_CLKRST\_LP\_BLETIMER\_DIV\_NUM}{12}{12}{{0}}%
\regfield{(reserved)}{12}{0}{000000000000}%
\reglabel{Reset}\regnewline%
\begin{regdesc}\begin{reglist}
\label{fielddesc:LPCLKRSTLPBLETIMERDIVNUM}\item [LP\_CLKRST\_LP\_BLETIMER\_DIV\_NUM] 配置 RTC BLE 定时器的时钟分频器的分频系数, RTC BLE 定时器的工作时钟频率= 源时钟频率／（( LP\_CLKRST\_LP\_BLETIMER\_DIV\_NUM－1）／2）
 (R/W)
\label{fielddesc:LPCLKRSTLPSELOSCSLOW}\item [LP\_CLKRST\_LP\_SEL\_OSC\_SLOW] 配置 RTC BLE 定时器的时钟源 \\
1: 选择  RC\_SLOW\_CLK 为 RTC BLE 定时器的源时钟 \\
0 : 禁止 RC\_SLOW\_CLK 为 RTC BLE 定时器的源时钟 \\ (R/W)
\label{fielddesc:LPCLKRSTLPSELOSCFAST}\item [LP\_CLKRST\_LP\_SEL\_OSC\_FAST]配置 RTC BLE 定时器的时钟源 \\
1: 选择  RC\_FAST\_CLK 为 RTC BLE 定时器的源时钟 \\
0: 禁止 RC\_FAST\_CLK 为 RTC BLE 定时器的源时钟 (R/W)
\label{fielddesc:LPCLKRSTLPSELXTAL}\item [LP\_CLKRST\_LP\_SEL\_XTAL]配置 RTC BLE 定时器的时钟源 \\
1: 选择 XTAL\_CLK 为 RTC BLE 定时器的源时钟 \\
0: 禁止 XTAL\_CLK 为 RTC BLE 定时器的源时钟 \\(R/W)
\label{fielddesc:LPCLKRSTLPSELXTAL32K}\item [LP\_CLKRST\_LP\_SEL\_XTAL32K]配置 RTC BLE 定时器的时钟源 \\
1: 选择 XTAL32K\_CLK 为 RTC BLE 定时器的源时钟 \\
0: 禁止 XTAL32K\_CLK 为 RTC BLE 定时器的源时钟\\ (R/W)
\end{reglist}\end{regdesc}
\end{register}


\begin{register}{H}{LP\_CLKRST\_XTAL32K\_REG}{0x002C}\label{regdesc:LPCLKRSTXTAL32KREG}
\regfield{LP\_CLKRST\_DAC\_XTAL32K}{3}{29}{{3}}%
\regfield{LP\_CLKRST\_DBUF\_XTAL32K}{1}{28}{{0}}%
\regfield{LP\_CLKRST\_DGM\_XTAL32K}{3}{25}{{3}}%
\regfield{LP\_CLKRST\_DRES\_XTAL32K}{3}{22}{{3}}%
\regfield{(reserved)}{22}{0}{0000000000000000000000}%
\reglabel{Reset}\regnewline%
\begin{regdesc}\begin{reglist}
\label{fielddesc:LPCLKRSTDRESXTAL32K}\item [LP\_CLKRST\_DRES\_XTAL32K] 配置 DRES。 (R/W)
\label{fielddesc:LPCLKRSTDGMXTAL32K}\item [LP\_CLKRST\_DGM\_XTAL32K] 配置 DGM。 (R/W)
\label{fielddesc:LPCLKRSTDBUFXTAL32K}\item [LP\_CLKRST\_DBUF\_XTAL32K] 配置 DBUF。 (R/W)
\label{fielddesc:LPCLKRSTDACXTAL32K}\item [LP\_CLKRST\_DAC\_XTAL32K] 配置 DAC。 (R/W)
\end{reglist}\end{regdesc}
\end{register}

\begin{register}{H}{LP\_CLKRST\_DATE\_REG}{0x03FC}\label{regdesc:LPCLKRSTDATEREG}
\regfield{LP\_CLKRST\_CLK\_EN}{1}{31}{{0}}%
\regfield{LP\_CLKRST\_CLKRST\_DATE}{31}{0}{{0x2207280}}%
\reglabel{Reset}\regnewline%
\begin{regdesc}\begin{reglist}
\label{fielddesc:LPCLKRSTCLKRSTDATE}\item [LP\_CLKRST\_CLKRST\_DATE] 版本控制寄存器。(R/W)
\label{fielddesc:LPCLKRSTCLKEN}\item [LP\_CLKRST\_CLK\_EN] 配置强制开启寄存器的配置的时钟。
0：无效值\\
1：强制开启\\
 (R/W)
\end{reglist}\end{regdesc}
\end{register}

\begin{register}{H}{LP\_AON\_SYS\_CFG\_REG}{0x0034}\label{regdesc:LPAONSYSCFGREG}
\regfield{LP\_AON\_HPSYS\_SW\_RESET}{1}{31}{{0}}%
\regfield{LP\_AON\_FORCE\_DOWNLOAD\_BOOT}{1}{30}{{0}}%
\regfield{(reserved)}{30}{0}{000000000000000000000000000000}%
\reglabel{Reset}\regnewline%
\begin{regdesc}\begin{reglist}
\label{fielddesc:LPAONFORCEDOWNLOADBOOT}\item [LP\_AON\_FORCE\_DOWNLOAD\_BOOT] 触发 CPU 复位，切换芯片启动模式。\\
0：无影响\\
1：如果 \hyperref[fielddesc:EFUSEDISFORCEDOWNLOAD]{EFUSE\_DIS\_FORCE\_DOWNLOAD} 为 0，则触发 CPU 复位，将芯片启动模式强制从 SPI Boot 模式切换至 Joint Download Boot 模式。\\
(R/W)
\label{fielddesc:LPAONHPSYSSWRESET}\item [LP\_AON\_HPSYS\_SW\_RESET] 配置系统的软件复位。\\
0：不复位\\
1：复位\\
(WT)
\end{reglist}\end{regdesc}
\end{register}


\begin{register}{H}{LP\_AON\_CPUCORE0\_CFG\_REG}{0x0038}\label{regdesc:LPAONCPUCORE0CFGREG}
\regfield{(reserved)}{1}{31}{{0}}%
\regfield{LP\_AON\_CPU\_CORE0\_STAT\_VECTOR\_SEL}{1}{30}{{1}}%
\regfield{(reserved)}{1}{29}{{0}}%
\regfield{LP\_AON\_CPU\_CORE0\_SW\_RESET}{1}{28}{{0}}%
\regfield{(reserved)}{28}{0}{0000000000000000000000000000}%
\reglabel{Reset}\regnewline%
\begin{regdesc}\begin{reglist}
\label{fielddesc:LPAONCPUCORE0SWRESET}\item [LP\_AON\_CPU\_CORE0\_SW\_RESET] 配置 CPU 的软件复位。\\
0：不复位\\
1：复位\\ (WT)
\label{fielddesc:LPAONCORE0STATVECTORSEL}\item [LP\_AON\_CPU\_CORE0\_STAT\_VECTOR\_SEL] 配置是否使 CPU 从 RTC 快速内存启动。 \\
0：不从 RTC 快速内存启动\\
1：从 RTC 快速内存启动\\
(R/W)
\end{reglist}\end{regdesc}
\end{register}


\begin{register}{H}{LPPERI\_CLK\_EN\_REG}{0x0000}\label{regdesc:LPPERICLKENREG}
\regfield{(reserved)}{1}{31}{{0}}%
\regfield{LPPERI\_EFUSE\_CK\_EN}{1}{30}{{1}}%
\regfield{(reserved)}{30}{0}{000000000000000000000000000000}%
\reglabel{Reset}\regnewline%
\begin{regdesc}\begin{reglist}
\label{fielddesc:LPPERIEFUSECKEN}\item [LPPERI\_EFUSE\_CK\_EN] 配置是否使能 eFuse 控制器模块的时钟门控。 \\
0：禁止时钟信号通过\\
1：使能时钟信号通过\\(R/W)
\end{reglist}\end{regdesc}
\end{register}


\begin{register}{H}{LPPERI\_RESET\_EN\_REG}{0x0004}\label{regdesc:LPPERIRESETENREG}
\regfield{(reserved)}{1}{31}{{0}}%
\regfield{LPPERI\_EFUSE\_RESET\_EN}{1}{30}{{0}}%
\regfield{(reserved)}{30}{0}{000000000000000000000000000000}%
\reglabel{Reset}\regnewline%
\begin{regdesc}\begin{reglist}
\label{fielddesc:LPPERIEFUSERESETEN}\item [LPPERI\_EFUSE\_RESET\_EN] 配置 eFuse 控制器模块的软件复位。\\
0：不复位\\
1：复位\\(R/W)
\end{reglist}\end{regdesc}
\end{register}
